\chapter{O que o satélite gratuito separa, e a que preço o comercial resolve}
\label{cap:satelite}

A seção~\ref{sec:fig2} fecha a verificação de geração distribuída por imagem
gratuita. Restam duas perguntas: se a rota funciona para geração centralizada,
onde o arranjo tem centenas de hectares, e quanto custaria comprar a resolução
que o sensor gratuito não entrega.

\section{Construção e operação de usina centralizada}
\label{sec:fig7}

\begin{ficha}
  \item[Janela]      2018 a 2026, uma cena por ano
  \item[Amostra]     34 conjuntos, 306 medições
  \item[Unidade]     conjunto por ano
  \item[Fonte]       Sentinel-2 L2A pelo catálogo STAC público; ANEEL/SIGA
  \item[Teste]       Wilcoxon pareado, duas comparações
\end{ficha}

Para cada conjunto mede-se a razão entre o brilho no vermelho dentro do arranjo
e num anel de controle afastado:

\[ \rho(t)=\frac{\overline{B_{\text{red}}}\big|_{d\le 400\,\mathrm{m}}}
{\overline{B_{\text{red}}}\big|_{1{,}5\le d\le 3\,\mathrm{km}}} \]

O anel controla iluminação, estação e umidade do solo, que afetam as duas
regiões da mesma forma. As séries são alinhadas pela data de entrada em operação
registrada no SIGA, e não pelo calendário.

\begin{figure*}[t]
\centering
\includegraphics[width=\linewidth]{../figures/fig7_deteccao_construcao.pdf}
\caption{\textbf{A construção de usinas fotovoltaicas é detectável e datável por
Sentinel-2; a operação não é.}
\textbf{a}, Razão entre o brilho no vermelho dentro do arranjo e num anel de
controle de 1,5 a 3 km, por usina em linha cinza e mediana entre usinas em linha
vermelha, com faixa entre quartis, alinhadas pela data de entrada em operação do
SIGA.
\textbf{b}, Sentinel-2 em cor natural sobre o complexo Sol do Cerrado, em escala
de refletância fixa entre as datas. Os círculos marcam as 17 usinas do SIGA.
\textbf{c}, Comparação pareada da razão de brilho entre fases, uma linha por
usina, com a mediana em destaque, e valor de $p$ de Wilcoxon pareado.
\textbf{d}, Diferença entre o ano de pico do brilho relativo e o ano de entrada
em operação. As barras em vermelho indicam erro de até um ano.}
\label{fig:construcao}
\end{figure*}

Da fase de base para a de construção a razão sobe de 1,01 para 1,30, um ganho de
34\% na mediana, em 24 de 30 usinas, com $p = 1{,}2 \times 10^{-5}$. Da base
para a operação a mediana vai de 0,95 para 1,02, um ganho de 6,6\%, com 11 de 18
usinas em alta contra 7 em queda e $p = 0{,}09$.

\conclusao{A obra tem assinatura espectral e data a entrada em operação. A fase
operacional não se separa do entorno a 10 m de resolução.}

A construção remove vegetação de centenas de hectares, e solo exposto tem
refletância alta no vermelho. Na fase operacional não há sinal em nenhuma
direção. Módulos são escuros, mas a 10 m o pixel mistura painel, corredor, via
de serviço e solo entre fileiras.

\campo{O que isso responde da proposta original} A hipótese de cruzar registro
com evidência de satélite é executável para geração centralizada, com uma
inversão. O satélite confirma a usina pela cicatriz da obra, e não pelo arranjo
em si. Como a cicatriz é evento transitório, a auditoria exige monitoramento
contínuo da série, e não observação do estado atual.

\campo{Uma correção de método dominou o resultado} A primeira execução usou a
coordenada do conjunto publicada pelo ONS, que corresponde à subestação
coletora, e resultou nula. Em 31 de 37 conjuntos a subestação fica a mais de
750 m do centroide do arranjo, com mediana de 1,9 km. A máscara amostrava
terreno fora do arranjo. Refeita com as coordenadas de usina do SIGA, a
diferença entre fases aparece.

\begin{objecao}
A qualidade da geolocalização de referência domina o resultado, e qualquer
replicação depende dessa camada. Uma cena por ano fixa a datação em resolução
anual, com acerto exato em 21\% dos casos e erro de até um ano em 62\%. Os
círculos de 400 m não são o polígono do arranjo, e delinear por segmentação
melhoraria o contraste. O pulso indica remoção de vegetação, que também ocorre
em lavoura, mineração e loteamento; a atribuição depende do cruzamento com o
registro. Nada disso se estende a telhados.
\end{objecao}

\campo{Posição na literatura} A razão dentro e fora do arranjo é uma versão
simplificada do \textit{Temporal Cluster Matching} \cite{robinson2021}, que data
construções comparando a relação espectral dentro e fora da pegada. Aquele
método foi aplicado a usinas solares em Karnataka. O Global Renewables Watch
\cite{grw2025} estima datas de construção em escala global com imagem de alta
resolução. O que esta seção acrescenta é a validação contra o registro de
outorga e o resultado negativo sobre a fase operacional.

\section{Resolução, cobertura e custo de imagem comercial}
\label{sec:aux}

A curva de detectabilidade da seção~\ref{sec:fig2} posiciona qualquer catálogo
de sensor comercial. Os 21 sensores de um fornecedor típico caem sobre ela em
três grupos.

\begin{figure*}[t]
\centering
\includegraphics[width=\linewidth]{../figures/aux_geopera_capacidade.pdf}
\caption{\textbf{Os sensores de um catálogo comercial posicionados contra a
curva de detectabilidade da geração distribuída, e o custo de aquisição por
área.}
\textbf{a}, Resolução e faixa de imageamento dos sensores do catálogo.
\textbf{b}, Fração da geração distribuída registrada que atinge o piso de quatro
pixels em cada resolução, com os sensores do catálogo sobrepostos.
\textbf{c}, Custo de aquisição por instalação verificada, no cenário com
coordenada por instalação e no cenário de varredura municipal.}
\label{fig:comercial}
\end{figure*}

A transição de detectabilidade ocorre entre 1 m e 5 m. Abaixo de 1 m, 99,5\% da
geração distribuída é segmentável, e os 18 sensores submétricos do catálogo
estão todos no platô. Acima, a fração cai para 86,6\% a 2 m, 7,5\% a 5,3 m e
1,6\% no Sentinel-2. Os dois sensores hiperespectrais do catálogo ficam acima do
joelho da curva.

A faixa de 2 m custa entre AUD 2 e 3 por km\textsuperscript{2} e entrega a mesma
resolução do CBERS-4A WPM, que é público. Entre 0,80 m e 2 m há apenas arquivo,
sem programação de aquisição.

\conclusao{O custo de verificação é determinado pela ausência de localização no
registro, e não pela resolução disponível no mercado.}

Com coordenada por instalação, um recorte de 0,01 km\textsuperscript{2} a 30 cm
custa entre AUD 0,18 e 0,28. Sem ela, a varredura de 1.000 km\textsuperscript{2}
de um município custa entre AUD 18 mil e 28 mil. São cinco ordens de grandeza
entre os dois cenários.

\campo{O que fica viável} Varrer entre 100 e 400 km\textsuperscript{2} da área
ocupada de dois ou três municípios rurais custa entre AUD 1.800 e 11.200, e
permite auditar a contagem agregada contra o registro municipal. A verificação
por instalação na classe rural permanece inviável.

As 361.622 instalações de pessoa jurídica têm CEP completo no registro,
geocodificável em nível de rua. Para esse segmento, com arranjos medianos de 66
e 92 m\textsuperscript{2} em comércio e indústria, o recorte dirigido é viável,
com custo da ordem de AUD 0,20 por instalação.

\begin{objecao}
Preços e sensores são os publicados pelo fornecedor, em dólar australiano e sem
tributo, e mudam sem aviso. A profundidade de arquivo não é publicada por
sensor, o que é decisivo porque todo desenho deste relatório é longitudinal. O
catálogo é apenas óptico, sem radar de abertura sintética, que seria o
complemento imune a nuvem.
\end{objecao}

\campo{Onde está a análise energética por satélite} A abordagem em uso
operacional emprega satélite meteorológico geoestacionário, com índice de nuvem
convertido em irradiância e depois em geração por modelo de desempenho, a 1 ou
2 km e passo de 10 minutos. Ela mede a nuvem, e não o arranjo. É a abordagem do
capítulo~\ref{cap:calibracao}.
